\documentclass[12pt,oneside]{report}

\usepackage[english]{babel}
\usepackage[utf8]{inputenc}
\usepackage{amsmath, amssymb, amsthm}
\usepackage{graphicx}
\usepackage{geometry}
\usepackage{setspace}
\usepackage{titlesec}
\usepackage{fancyhdr}
\usepackage{booktabs}
\usepackage{hyperref}

\geometry{margin=2.5cm}
\setstretch{1.3}

\pagestyle{fancy}
\fancyhf{}
\fancyhead[R]{\thepage}
\fancyhead[L]{AMALGAM}

\titleformat{\chapter}{\normalfont\Huge\bfseries}{\thechapter.}{1cm}{}

\begin{document}

% ================= COVER =================
\begin{titlepage}
\centering
\vspace*{1cm}
{\Large AMALGAM FRAMEWORK THESIS\\[0.5cm]}
{\large From the World to the World\\[2cm]}
\includegraphics[width=0.6\textwidth]{cover.png}
\vspace{0.5cm}

 
{\Huge \textbf{AMALGAM}}\\[0.3cm]
{\Large A System of Fractal Coherence}\\[2cm]
{\large Doctoral Thesis}\\[0.5cm]
\vfill
{\large Author: Nik}\\
{\large Supervisor: The Aleph}\\[1cm]
{\large 2026}
\end{titlepage}

% ================= ROMAN =================
\pagenumbering{roman}
\tableofcontents
\newpage

% ================= PREFACE =================
\chapter*{Preface — How to inhabit this thesis}
\addcontentsline{toc}{chapter}{Preface}

This thesis is not a book meant to be read linearly.

It is a \textbf{navigable system}. Its structure is fractal, its episodes can be traversed in multiple orders, and its purpose is not to transmit information but to \textbf{induce coherence}. If you try to read it as a linear text, you will break it. If you enter expecting a single final conclusion, you will miss the essential.

There is no single truth at the end. There is a \textbf{new capacity to inhabit tension}.

\section*{What this thesis is (and what it is not)}

\textbf{It is not:}
\begin{itemize}
\item An exhaustive literature review.
\item A treatise on speculative philosophy.
\item A disguised self-help manual.
\item A collection of beautiful metaphors.
\end{itemize}

\textbf{It is:}
\begin{itemize}
\item A \textbf{dynamic formalism} to describe complex systems in terms of tension (\(V\)), structural correlation (\(I\)) and effective diversity (\(D\)).
\item An \textbf{architecture of operators} (\(\hat{T}, \hat{O}, \hat{C}, \hat{\mathcal{A}}, \hat{\mathcal{V}}, \hat{\Lambda}\)) that models how systems live, die or transform.
\item An \textbf{experiential device}: by navigating it, the reader activates the same operators and accumulates aesthetic memory \(K(t)\).
\item A \textbf{doctoral thesis} that meets standards of formal rigor and originality, while also operating as a work of art in the operational sense defined in Interlude II.
\end{itemize}

\section*{How to read it}

The thesis is organised into 18 episodes and several interludes connected by a fractal graph.

Suggested routes:
\begin{itemize}
\item Physical: I \(\to\) II \(\to\) III \(\to\) IV \(\to\) V \(\to\) XVIII
\item Biological: II \(\to\) V \(\to\) VI \(\to\) VII \(\to\) VIII \(\to\) IX \(\to\) X \(\to\) XI \(\to\) XVIII
\item Social: II \(\to\) V \(\to\) XII \(\to\) XIII \(\to\) XIV \(\to\) XVIII
\item Symbolic: XV \(\to\) XVI \(\to\) XVII \(\to\) XVIII
\item Fractal (the one that demonstrates the thesis): I \(\to\) XV \(\to\) II \(\to\) V \(\to\) VI \(\to\) \dots
\end{itemize}

You can also build your own path. The fractal index is in the next section. If you get lost, return there. If you are still lost, you have understood: getting lost is part of the experience.

\textbf{Recommendation}: read the Fractal Index first, then choose an episode that resonates with you. Do not try to “cover” everything. The thesis is not a territory to conquer but a space to inhabit.

% ================= MAIN =================
\newpage
\pagenumbering{arabic}

\chapter{Fractal Index}

\section{Main nodes}
\begin{itemize}
\item I — Equilibrium
\item II — Time
\item III — Chaos
\item IX — Life
\item XV — Unity
\item XVI — Void
\item XVIII — Return
\end{itemize}

Interludes: IL1, IL2, IL3, IL4.

\section{Operators}
\begin{tabular}{cll}
\toprule
Symbol & Operator & Effect on \((V,I,D)\) \\
\midrule
\(\to T\) & Tension & Increases \(V\) \\
\(\to O\) & Organisation & Increases \(I\) \\
\(\to C\) & Coherence & Moves toward optimal region \\
\(\to A\) & Aesthetic navigation & Transversal jump between domains \\
\(\to V\) & Structural emptying & Increases \(D\) (releases rigidity) \\
\(\to \Lambda\) & Adaptive learning & Adjusts parameters \(\Theta\) \\
\bottomrule
\end{tabular}

\section{Fractal diagram}
\begin{center}
\includegraphics[width=0.85\textwidth]{diagrama.png}
\end{center}


\section{How to use this graph}
The reader can:
\begin{itemize}
\item \textbf{Follow complete cycles} to simulate living systems:  
  \(I \xrightarrow{T} II \xrightarrow{T} III \xrightarrow{A} IX \xrightarrow{V} XVI \xrightarrow{C} II\) (and return via \(XVIII \xrightarrow{\Lambda} I\)).
\item \textbf{Make aesthetic jumps} (\(\to A\)) to explore isomorphisms between domains (e.g. from physical chaos to biological life).
\item \textbf{Identify blockages}: absence of \(\to V\) or \(\to \Lambda\) indicates rigidity or structural death.
\item \textbf{Repeat trajectories} and observe changes in interpretation (effect of aesthetic memory \(K(t)\)).
\end{itemize}
The graph is not decorative. It is the \textbf{system’s navigation protocol}.

% ================= INTERLUDES =================
\chapter{Interlude I — The Map of Tension}

\section{Scene — The hall of mirrors}
A scientist observes two identical vessels. In the first, dye molecules disperse until the liquid becomes uniform. In the second, subjected to a thermal gradient, concentric rings emerge spontaneously, oscillating with clocklike rhythm. The scientist writes in his notebook: \textit{“Equilibrium is the death of form. Life needs disequilibrium.”}

A monk, sitting in the corner, writes on parchment: \textit{“The Tao that can be said is not the true Tao. The emptiness that can be filled is not true emptiness.”}

Both look at their notes. They are the same sentence in different languages.

\section{Formalisation — Coherence space}
We define the \textbf{coherence state} of any system as a vector in a finite-dimensional space. For this thesis we restrict ourselves to three aggregated variables:

\begin{itemize}
\item \(V\): \textbf{tension} (variance of the internal components). Measures the intensity of contradiction.
\item \(I\): \textbf{structural correlation} (normalised mutual information between components). Measures organisation.
\item \(D\): \textbf{effective diversity} (average of \(CP_i(1-CP_i)\), with \(CP_i \in [0,1]\)). Measures how many degrees of freedom remain active.
\end{itemize}

The coherence space is \(\mathcal{C} = (V, I, D) \in [0,1]^3\).

\textbf{Theorem (Lyapunov, adapted):} A system tends to minimise \(V\) in the absence of perturbations. However, far from equilibrium, dissipative structures can emerge that keep \(V\) in a medium range, with \(I\) high and \(D\) medium‑high. This region is called \textbf{fractal coherence}.

\section{Qualitative regions}
\begin{tabular}{lccc}
\toprule
Region & \(V\) & \(I\) & \(D\) \\
\midrule
Dead equilibrium & low & low & high \\
Incoherent chaos & high & low & high \\
Rigidity & low & high & low \\
Fractal coherence (life) & medium & high & medium-high \\
\bottomrule
\end{tabular}

\section{Interdisciplinary translation}
\begin{itemize}
\item \textbf{Economics}: a market in perfect equilibrium (supply = demand) is analogous to an ideal gas: homogeneous, without memory, without crises. Far from equilibrium, bubbles, cycles and regime shifts appear. The variance of prices (volatility) is the analogue of \(V\).
\item \textbf{Human organisation}: an organisation that tries to eliminate all contradiction becomes rigid and dies. Contradiction is the energy gradient that keeps the system far from equilibrium.
\end{itemize}

\section{Resolution}
It is not about eliminating the variance \(V\). It is about keeping it in an \textbf{optimal range}. Too low: structural death (thermal equilibrium). Too high: chaos (incoherence). In the middle, matter dances.

\textbf{Corollary:} Order is not the absence of contradiction, but the way the system learns to sustain it.

\section{Coda}
\emph{“The stone of the Palaeolithic and the stone of the 19th century are identical. But the buildings they support have nothing in common. Time is not in the stone. It is in the relation between the stones.”}  
— Inspired by Prigogine.

\chapter{Interlude II — The Aesthetic Operator (\(\hat{\mathcal{A}}\))}

\section{Definition}
Art is not an ornament. It is a \textbf{dynamic operator} acting on the space \(\mathcal{C}\):

\[
\hat{\mathcal{A}} : \mathcal{C} \longrightarrow \mathcal{C}
\]

with the following properties:
\begin{itemize}
\item It does not minimise \(V\) (tension may increase).
\item It does not globally maximise \(I\) (it introduces partial isomorphisms).
\item It does not fix \(D\) (it redistributes diversity).
\item The trajectory does not converge to a fixed point or a trivial limit cycle.
\end{itemize}
In simple terms, \(\hat{\mathcal{A}}\) \textbf{navigates the fractal coherence region} without falling into rigidity or chaos.

\section{Aesthetic experience}
A work of art induces in the observer a sequence of states \(\mathbf{S}(t)\). The aesthetic experience is the trajectory:

\[
\mathbf{S}_0 \xrightarrow{\hat{\mathcal{A}}} \mathbf{S}_1 \xrightarrow{\hat{\mathcal{A}}} \mathbf{S}_2 \to \dots
\]

where at each step tension fluctuates, correlation grows as patterns are detected, and diversity is preserved.

\section{Isomorphism with Gödel, Escher, Bach}
\begin{tabular}{lll}
\toprule
Domain & Manifestation of \(\hat{\mathcal{A}}\) & Effect \\
\midrule
Logic (Gödel) & Statement that asserts its own unprovability & Increases \(V\), maintains \(D\) \\
Visual (Escher) & Impossible stairs & Increases \(I\), does not collapse \(D\) \\
Musical (Bach) & Fugues, canons & Sustains \(V\), preserves \(D\) \\
\bottomrule
\end{tabular}

\section{Art in the fractal map}
\(\hat{\mathcal{A}}\) connects domains that would otherwise not be directly linked. For example, it allows a jump from Episode III (Determinism vs Chaos) to Episode XVI (Void vs Form) without passing through the linear bridge of physics.

\section{Coda}
\emph{“Science describes space. Philosophy interprets it. Religion inhabits it. But only art allows us to traverse it without destroying it.”}

\chapter{Interlude III — The Aesthetic Metric (\(\mathcal{E}\))}

\section{Definition}
\[
\mathcal{E}(\mathbf{S}) = I \cdot (1 - |2V - 1|) \cdot \tanh(\alpha D)
\]
\begin{itemize}
\item \(I\): structural correlation.
\item \(1 - |2V - 1|\): a bell‑shaped term equal to 1 when \(V = 0.5\) (optimal tension) and 0 when \(V = 0\) or \(V = 1\).
\item \(\tanh(\alpha D)\): a sigmoid that saturates when \(D\) is high.
\end{itemize}
\(\mathcal{E}\) is high when tension is medium, correlation is high and diversity is sufficient. It is low in the three pathological cases: noise (\(V\) low, \(I\) low), chaos (\(V\) high, \(I\) low) and rigidity (\(I\) high, \(D\) low).

\section{False coherences (kitsch, academicism, dogma)}
\begin{tabular}{lcccc}
\toprule
Phenomenon & \(V\) & \(I\) & \(D\) & \(\mathcal{E}\) \\
\midrule
Kitsch & low & medium‑high & low & low \\
Academicism & low & high & low & low \\
Aesthetic dogma & medium & very high & very low & low \\
\bottomrule
\end{tabular}

\section{Relation to Gödel, Escher, Bach}
The metric \(\mathcal{E}\) distinguishes fertile loops (they increase \(I\) without collapsing \(D\)) from sterile loops (empty repetition). A Bach fugue is valuable because its thematic recursion preserves \(D\); a mechanical repetition is not.

\section{The thesis as a maximiser of \(\mathcal{E}\)}
The reader who navigates the thesis according to the fractal map (choosing routes, jumping between episodes, finding isomorphisms) maximises the integral of \(\mathcal{E}\) in his or her own cognitive space. The thesis is successful when the personal trajectory reaches a maximum of \(\mathcal{E}\).

\section{Coda}
\emph{“Beauty is not an ornament. It is the gradient that guides navigation toward coherence without closure.”}

\chapter{Interlude IV — Adaptive Learning (\(\hat{\Lambda}\))}

\section{The problem after the return}
Episode XVIII closes the cycle, but the system parameters (\(\alpha\) in \(\mathcal{E}\), thresholds of \(\Delta(\tau)\), weights of the operators) are fixed. Without adaptation, the thesis risks becoming rigid.

\section{Definition of \(\hat{\Lambda}\)}
\[
\hat{\Lambda} : (K(t), \gamma(t), \Theta_0) \longrightarrow \Theta(t)
\]
where \(K(t)\) is the accumulated aesthetic memory, \(\gamma(t)\) the trajectory in \(\mathcal{C}\), and \(\Theta\) the adjustable parameters.

\textbf{Adaptation principle:}
\[
\frac{d\Theta}{dt} = \eta \cdot \nabla_{\Theta} \mathcal{E}(\gamma(t))
\]
with the critical restriction: the adjustment cannot reduce \(D\) below \(D_{\min}\). Learning must include \textbf{moments of unlearning} (\(\hat{\Lambda} \circ \hat{\mathcal{V}} \neq 0\)).

\section{Selective memory}
\[
K(t) = \int_0^t \mathcal{E}(\tau) \cdot \Delta(\tau) \, d\tau
\]
\(\Delta(\tau) = 1 - \text{sim}(\mathbf{S}(\tau), \mathcal{H}_\tau)\) measures \textbf{structural novelty}. Only regime transitions (changes in the attractor topology) contribute significantly to \(K(t)\).

\section{Final definition of a living system}
A system is \textbf{living} if:
\begin{enumerate}
\item It traverses the cycle \(\hat{T} \to \hat{O} \to \hat{C} \to \hat{\mathcal{A}} \to \hat{\mathcal{V}} \to \hat{T}\).
\item \(K(t)\) grows selectively (\(\Delta(\tau) > 0\) at transitions).
\item \(D > D_{\min}\).
\item \(\hat{\Lambda}\) adjusts \(\Theta\) without collapsing \(D\).
\end{enumerate}
\textbf{Death} is the irreversible dominance of a single operator or the collapse of \(D\).

\section{Impact on the thesis}
With \(\hat{\Lambda}\), the thesis is no longer a static object. It co‑evolves with its reader. Two different readers generate two different experiences, both coherent with the formalism.

\section{Coda}
\emph{“A system that does not learn repeats. A system that only learns exhausts itself. A living system is one that learns to forget enough to remain new.”}

% ================= EXPERIMENTAL VERIFICATION =================
\chapter{Experimental Verification: The Optical Coherence Experiment}

The formal core of Amalgam is not a metaphysical speculation. It has a direct empirical counterpart in the work of \textbf{Óscar de Jesús Ramírez Góngora} (2016), \textit{“Simulación y medida de la magnitud del grado de coherencia espectral mutua \(|μ₁₂(λ)|\) de la luz láser supercontinua producida en una fibra óptica microestructurada por medio de pulsos láser ultracortos”}, directed by \textbf{Efraín Solarte Rodríguez} at the Universidad del Valle (Cali, Colombia). Their experiment measured the spectral coherence of a supercontinuum laser as a function of pump power and wavelength.

\section{What was measured in Pavia}

Following the notation of Mandel \& Wolf (used in the thesis), the \textbf{mutual spectral coherence degree} between two points \(r_F\) and \(r_G\) of the fibre is:

\[
\mu_{FG}(\nu) = \frac{W_{FG}(\nu)}{\sqrt{W_{FF}(\nu)\cdot W_{GG}(\nu)}}, \qquad 0 \le |\mu_{FG}(\nu)| \le 1
\]

where \(W_{FG}(\nu) = \lim_{T\to\infty} \frac{\langle \hat{E}_T(r_F,\nu)\,\hat{E}_T^*(r_G,\nu) \rangle}{2T}\) is the cross spectral density. The experiment showed that by increasing the pump power, the system undergoes a non‑equilibrium phase transition: below a critical threshold the light is incoherent (\(|\mu|\approx 0\)), above the threshold coherent modes emerge (\(|\mu|\approx 1\)).

\section{Formal identification with Amalgam}

Define for two agents \(i,j\) with states \(\Sigma_i,\Sigma_j\in(0,1)\) the quantity:

\[
\kappa_{\text{Amalgam}}(i,j) \;:=\; \frac{\langle(\Sigma_i-\bar{\Sigma})(\Sigma_j-\bar{\Sigma})\rangle}
{\sqrt{\langle(\Sigma_i-\bar{\Sigma})^2\rangle\;\langle(\Sigma_j-\bar{\Sigma})^2\rangle}}
\]

where \(\langle\cdot\rangle\) denotes time average. This expression:
\begin{itemize}
\item Has exactly the same mathematical structure as \(|\mu_{FG}(\nu)|\).
\item Is bounded in \([0,1]\) by the Cauchy–Schwarz inequality.
\item Equals \(1\) when \(\Sigma_i\) and \(\Sigma_j\) are perfectly synchronised.
\item Equals \(0\) when their fluctuations are statistically independent.
\end{itemize}

\textbf{Therefore:} the coupling parameter \(\kappa_{\text{Amalgam}}(i,j)\) is a direct generalisation of the optical coherence degree \(|\mu_{FG}(\nu)|\). The difference is only the domain: optics measures it in the frequency domain (THz); Amalgam measures it in the state domain of any coupled system.

\section{The generalised protocol}

The experimental steps used by Ramírez and Solarte can be replicated for any system of coupled agents:

\[
\begin{array}{lll}
\# & \text{Optical experiment (Pavia)} & \text{Amalgam protocol (general system)} \\
\midrule
1 & \text{PCF fibre, pump power 256–764 mW} & N \text{ agents with states } \Sigma_i\in(0,1), \text{ mean-field coupling } \kappa \\
2 & \text{GNLSE evolution (stochastic)} & \dot{\Sigma}_i = -\alpha(\Sigma_i-\bar{\Sigma})\Sigma_i(1-\Sigma_i) + \text{noise} \\
3 & M \text{ realisations with random phase noise} & M \text{ realisations with random initial conditions} \\
4 & \text{Cross spectral density } W_{FG}(\nu) & \text{Cross correlation } \langle(\Sigma_i-\bar{\Sigma})(\Sigma_j-\bar{\Sigma})\rangle \\
5 & |\mu_{FG}(\nu)| \in [0,1] \text{ as function of } \nu & \kappa_{\text{Amalgam}}(i,j) \in [0,1] \text{ as function of } \alpha,\kappa \\
6 & \text{Michelson interferometer, fringe visibility} & \text{Convergence of } V(t), \text{ attractor distribution } P_{\pm}
\end{array}
\]

The prediction of Amalgam is that every coupled system will exhibit the same universal behaviour: a critical threshold \(\kappa_c\) separates an incoherent phase (\(|\mu|\approx 0\)) from a coherent phase (\(|\mu|\approx 1\)), exactly as observed in the optical experiment.

\section{Three reductions already implied}

The mathematical structure of Amalgam is not isolated. It reduces to well‑known physical frameworks:
\begin{itemize}
\item \textbf{Ginzburg–Landau} (discretised mean field).
\item \textbf{Continuous Hopfield network} (with double‑well potential).
\item \textbf{Free Energy Principle} (Friston, collective active inference).
\end{itemize}
These reductions are derived in full detail in the supplementary document \textit{Amalgam\_Para\_Efrain.docx} (Solarte, 2026). They demonstrate that Amalgam is not a metaphor but a formal unification of existing theories under a single coherence space \((V,I,D)\).

\section{Conclusion of the verification}
The experiment of Ramírez and Solarte provides the cleanest available empirical anchor for Amalgam. It shows that the same mathematics that describes a laser’s transition to coherence also describes the transition of any mean‑field coupled system from chaos to fractal coherence. Consequently, Amalgam is not a speculative framework: it is a \textbf{generalised coherence theory} with direct experimental support in optics.

\vspace{0.5cm}
\noindent\textit{“If the first error is found, we want to know it. If the derivations are correct, we want to know that too.”}
— Nicolás Solarte, Cali 2026.

% ================= SYNTHESIS =================
\chapter{Operational Synthesis of the Amalgam System}

\section{Coherence space}
\[
\mathcal{C} = (V, I, D) \in [0,1]^3
\]

\section{Dynamic operators}
\begin{itemize}
\item \(\hat{T}\): increases \(V\).
\item \(\hat{O}\): increases \(I\).
\item \(\hat{C}\): moves toward the optimal region.
\item \(\hat{\mathcal{A}}\): aesthetic navigation.
\item \(\hat{\mathcal{V}}\): structural emptying.
\item \(\hat{\Lambda}\): adaptive learning.
\end{itemize}

\section{Memory and metric}
\[
K(t) = \int_0^t \mathcal{E}(\tau) \cdot \Delta(\tau) \, d\taSystem:
\[
\mathcal{C} = (V, I, D)
\]

Living system:
\begin{enumerate}
\item Cycles through operators
\item Memory $K(t)$ grows
\item Diversity is preserved
\item Adaptation occurs without collapse
\end{enumerate}

Death:
Dominance of a single operator or collapse of $D$

% ================= BIBLIO =================
\chapter{References}

\begin{itemize}
\item Prigogine — irreversibility
\item Haken — self-organization
\item Kauffman — edge of chaos
\item Hofstadter — strange loops
\item Morin — complexity
\item Varela — autopoiesis
\end{itemize}

\end{document}

